\documentclass[11pt]{article}

\usepackage[a4paper,margin=1in]{geometry}
\usepackage{amsmath,amssymb,amsthm,mathtools}
\usepackage{enumitem}
\usepackage[hidelinks]{hyperref}

\newtheorem{theorem}{Theorem}
\newtheorem{proposition}[theorem]{Proposition}
\newtheorem{lemma}[theorem]{Lemma}
\newtheorem{corollary}[theorem]{Corollary}
\theoremstyle{definition}
\newtheorem{definition}[theorem]{Definition}
\newtheorem{remark}[theorem]{Remark}
\newtheorem{fact}[theorem]{Fact}

\newcommand{\R}{\mathbb R}
\newcommand{\Om}{\Omega}
\newcommand{\HH}{\mathcal H}
\newcommand{\eps}{\varepsilon}
\newcommand{\dT}{\delta_T}
\newcommand{\diso}{\delta_{\mathrm{iso}}}
\newcommand{\AF}{\mathcal A_F}
\newcommand{\gd}{\mathrm{gd}}
\newcommand{\Etil}{\widetilde E}
\newcommand{\Exc}{\mathrm{Exc}}
\newcommand{\norm}[1]{\left\|#1\right\|}
\newcommand{\abs}[1]{\left|#1\right|}
\newcommand{\dstar}{\delta_{*}}

\title{Sharp planar Saint--Venant stability\\
by reduction to a graph inset\\
\large selection-free, transport-free: only $D_H$ and $D_I$}
\author{}
\date{}

\begin{document}
\maketitle

\begin{abstract}
We prove the sharp planar quantitative Saint--Venant (Faber--Krahn for the
torsion) inequality $\AF(\Om)^2\le C\,\dT(\Om)$ by reducing $\Om$ to a
small-norm radial graph $G$ with $\abs{\Om\triangle G}\lesssim\sqrt\dT$ and
$\dstar(G)\lesssim\dT$, and then invoking the perturbative radial-graph
theorem (G). \emph{The reduction uses nothing beyond the two level-set
defects} $D_H$ (the Cauchy--Schwarz/Serrin defect) and $D_I$ (the
isoperimetric defect): no selection principle, no moving-ball or kinetic
transport, no Magnanini--Poggesi stability. The two ingredients that make
the reduction \emph{sharp} are (i) a fold-content bound with no
quantitative-isoperimetric $\sqrt{\,\cdot\,}$ loss, obtained from the
point-source flux identity together with the cancellation
$\abs{I_{\Etil}-2\pi\hat r}\le C a\eta$ across the thin trapping annulus; and
(ii) the observation that cutting the (low-torsion) caps to manufacture a
graph \emph{decreases} the scale-invariant torsion deficit, because the
equal-area ball shrinks faster than the torsion is lost. Choosing the inset
level at the scale $\hat v\sim\sqrt\dT$ then balances closeness against the
graph defect at exactly $\sqrt\dT$.
\end{abstract}

\tableofcontents

\section{Setup, conventions, and the two defects}
\label{sec:setup}

Let $\Om\subset\R^2$ be open with $\abs\Om=\pi$, and let $u\in
H^1_0(\Om)$ be its torsion function, $-\Delta u=1$ in $\Om$, $u=0$ on
$\partial\Om$. Write $T(\Om)=\int_\Om u$ for the torsional rigidity and
\[
   \dT:=\dT(\Om):=T(B_1)-T(\Om)=\tfrac\pi8-T(\Om)\ \ge\ 0
\]
for the (absolute, area-normalised) Saint--Venant deficit; $B_1$ is the unit
disc, $T(B_1)=\pi/8$. For a measurable $E$ we write $\AF(E)$ for the
Fraenkel asymmetry,
\[
   \AF(E):=\min_{x_0}\frac{\abs{E\triangle B(x_0,r_E)}}{\abs E},\qquad
   r_E=\sqrt{\abs E/\pi},
\]
and, for the inset arguments below, the \emph{scale-invariant} relative
torsion deficit
\begin{equation}\label{eq:dstar}
   \dstar(E):=\frac{T(B_E)-T(E)}{T(B_E)}\in[0,1),\qquad
   B_E:=\text{any ball with }\abs{B_E}=\abs E,\quad T(B_E)=\frac{\abs E^2}{8\pi}.
\end{equation}
Both $\AF$ and $\dstar$ are invariant under translation and dilation. For
$\Om$ itself, $\dstar(\Om)=(\tfrac\pi8-T(\Om))/(\pi/8)=\tfrac8\pi\dT$, so
$\dstar(\Om)\asymp\dT$; the target $\AF(\Om)^2\le C\dT$ and
$\AF(\Om)^2\le C'\dstar(\Om)$ are equivalent.

\paragraph{Distribution function and the two defects.}
For $v\in(0,m)$, $m:=\max u$, set
\[
   \mu(v):=\abs{\{u>v\}},\qquad P(v):=P(\{u>v\})=\HH^1(\{u=v\})\ \text{(a.e. }v),
\]
and define the \emph{Serrin/Cauchy--Schwarz defect} and the
\emph{isoperimetric defect} of the level $v$,
\begin{equation}\label{eq:DHDI}
   D_H(v):=\mu(v)\bigl(-\mu'(v)\bigr)-P(v)^2\ \ge0,
   \qquad
   D_I(v):=P(v)^2-4\pi\mu(v)\ \ge0 .
\end{equation}
Here $D_H\ge0$ is Cauchy--Schwarz applied to
$\int_{\{u=v\}}\!1=\int|\nabla u|^{1/2}|\nabla u|^{-1/2}$ against the flux
and coarea identities
\begin{equation}\label{eq:fluxcoarea}
   \int_{\{u=v\}}\!\abs{\nabla u}\,d\HH^1=\mu(v),
   \qquad
   -\mu'(v)=\int_{\{u=v\}}\!\frac{d\HH^1}{\abs{\nabla u}},
\end{equation}
and $D_I\ge0$ is the isoperimetric inequality for $\{u>v\}$; $D_H(v)=0$ iff
$\abs{\nabla u}$ is constant on $\{u=v\}$, and $D_I(v)=0$ iff $\{u>v\}$ is a
disc. Everything in this note is built from $D_H$ and $D_I$.

\section{Imported facts}
\label{sec:imports}

We collect the inputs proved elsewhere in this project or standard in the
literature. Each is used as a black box; none uses selection, transport, or
Magnanini--Poggesi.

\begin{fact}[Level identity; \texttt{sec\_core.tex}, Thm.~``identity'']
\label{fact:identity}
With the normalisation $\abs\Om=\pi$,
\begin{equation}\label{eq:identity}
   \int_0^m\bigl(D_H(v)+D_I(v)\bigr)\,dv\ =\ 4\pi\,\dT .
\end{equation}
The identity is dimensionally homogeneous of degree four and hence
scale-invariant: applied to any domain with its own torsion function it
reads $\int(D_H+D_I)=4\pi\,(T(\text{ball})-T(\text{domain}))$.
\end{fact}

\begin{fact}[Talenti comparison and the profile identity;
\texttt{sec\_core.tex}, Thm.~``talenti'']\label{fact:talenti}
Let $\mu_B(v):=\pi(1-4v)_+$ be the distribution function of the disc
torsion. Then $\mu(v)\le\mu_B(v)$ for all $v$, and
\begin{equation}\label{eq:profile}
   \int_0^{1/4}\bigl(\mu_B(v)-\mu(v)\bigr)\,dv\ =\ \dT,
   \qquad \mu_B-\mu\ \ge0 .
\end{equation}
In particular $m\le1/4$. Moreover, from $D_H,D_I\ge0$ and
\eqref{eq:fluxcoarea}, $-\mu'(v)\ge4\pi$ for a.e.\ $v$.
\end{fact}

\noindent
(For \eqref{eq:profile}: $T(B_1)=\int_0^{1/4}\mu_B$, $T(\Om)=\int_0^m\mu$,
and $\mu\le\mu_B$ with $\mu\equiv0$ on $(m,\tfrac14]$; for $-\mu'\ge4\pi$:
$\mu(-\mu')\ge P^2\ge4\pi\mu$ by \eqref{eq:DHDI}.)

\begin{fact}[Inset regularity; \texttt{sec\_smoothness.tex},
\texttt{sec\_core.tex}]\label{fact:reg}
For a.e.\ \emph{regular} value $v$ of $u$ the super-level set $\{u>v\}$ has
finite perimeter with $\HH^1$-rectifiable, piecewise-analytic boundary; the
flux and coarea identities \eqref{eq:fluxcoarea} hold; and (superharmonicity
of $u$, minimum principle) every connected component of $\{u<v\}$ touches
$\partial\Om$, so a connected component of $\{u>v\}$ has no interior holes.
By Sard's theorem the non-regular values have measure zero.
\end{fact}

\begin{fact}[Trapping annulus; \texttt{sec\_density\_annulus.tex},
Thm.~``perim-annulus'']\label{fact:annulus}
Let $E$ be a connected set of finite perimeter without interior holes,
$\abs E=\mu_E$, $\hat r=\sqrt{\mu_E/\pi}$, with isoperimetric deficit
$\diso(E)$ and perimeter excess $\Exc:=P(E)-2\pi\hat r=2\pi\hat r\,\diso(E)$.
Let $a:=\abs{E\triangle B_{\hat r}(z)}$ for the optimal centre $z$. Then $E$
is trapped in a concentric annulus,
\begin{equation}\label{eq:annulus}
   B_{\rho_i}(z)\subset E\subset B_{\rho_e}(z),\qquad
   \eta:=\rho_e-\rho_i\ \le\ C\,(\Exc+a)^{1/2},
\end{equation}
with $\rho_i\le\hat r\le\rho_e$ and an absolute constant $C$. (Proof: the
angular-perimeter coarea identity
$\int_{\partial E}\abs{\nu_\theta}=\int_0^\infty N(\rho)\,d\rho$ with
$N\ge2$ on $(\rho_i,\rho_e)$, Cauchy--Schwarz, and the radial-profile bound;
no $C^2$ geometry, no covering.)
\end{fact}

\begin{fact}[Quantitative isoperimetry, planar Fuglede/FMP]\label{fact:fmp}
There is an absolute $C_{\mathrm F}$ with
$\AF(E)\le C_{\mathrm F}\sqrt{\diso(E)}$, i.e.
$a=\abs{E\triangle B_{\hat r}(z)}\le C_{\mathrm F}\sqrt{\diso(E)}\,\abs E$.
\end{fact}

\begin{fact}[Point-source flux identities; \texttt{sec\_graph.tex},
Lem.~``two-flux'']\label{fact:flux}
Let $E$ be of finite perimeter with $z$ in its interior, $e_r=(x-z)/|x-z|$,
$\nu_r=\nu\cdot e_r$ the radial part of the outward normal, and for
$\theta\in S^1$ let $M(\theta)=\HH^0(\partial E\cap\{z+\rho e^{i\theta}:
\rho>0\})$ be the number of crossings of the ray. Then
\begin{equation}\label{eq:two-flux}
   \int_{\partial E}\frac{\abs{\nu_r}}{\abs{x-z}}\,d\HH^1
   =\int_{S^1}M(\theta)\,d\theta,
   \qquad
   \int_{\partial E}\frac{\nu_r}{\abs{x-z}}\,d\HH^1=2\pi,
   \qquad
   \int_{\partial E}\nu_r\,d\HH^1=\int_E\frac{dx}{\abs{x-z}}=:I_E .
\end{equation}
\end{fact}

\begin{fact}[Boundary-layer torsion stability;
\texttt{selection\_free\_route.tex}, Lem.~``torsion-stab'']%
\label{fact:torsion-stab}
Let $E,F$ with $B_{\rho_i}(z)\subset E\cap F$ and $E\cup F\subset
B_{\rho_e}(z)$, $\abs{E\triangle F}=m$, so $E\triangle F$ lies in the collar
$\{\rho_i<\abs{x-z}<\rho_e\}$. Then
$\abs{T(E)-T(F)}\le 3\,\tilde\eta\,m$, where
$\tilde\eta:=\tfrac14(\rho_e^2-\rho_i^2)$ bounds the torsion of $B_{\rho_e}$
on the collar. (Green-function double integral, monotonicity
$G_{E\cap F}\le G_{E\cup F}$, and the barrier $u_{E\cup F}\le\tilde\eta$ on
the collar.)
\end{fact}

\begin{fact}[Radial-graph theorem (G);
\texttt{selection\_free\_route.tex}, Module (G)]\label{fact:graph}
There are absolute constants $\eta_G=\tfrac1{13}$ and $C_G<48$ such that: if
$G=\{z+r\omega:0\le r<\rho(\omega)\}$ with $\rho\in C^0(S^1)$,
$\norm{\rho/\hat r_G-1}_{L^\infty}\le\eta_G$ \textup{(}$\hat r_G=
\sqrt{\abs G/\pi}$\textup{)} and barycentre $z$, then
\begin{equation}\label{eq:graphG}
   \AF(G)^2\ \le\ C_G\,\dstar(G).
\end{equation}
The statement is scale-invariant; an admissible $G$ need not have area $\pi$.
It uses no property of $\Om$ and no selection.
\end{fact}

\section{The inset, its deficit, and its closeness to $\Om$}
\label{sec:inset}

Throughout, $\delta_\star>0$ denotes a small absolute constant, redefined
finitely many times, and $\dT\le\delta_\star$.

\subsection{Deficit transfer to any super-level set}

\begin{lemma}[Deficit transfer]\label{lem:transfer}
For every regular $v\in(0,m)$ the connected components of $\{u>v\}$ each
have scale-invariant deficit controlled by $\Om$. Precisely, writing
$E_v:=\{u>v\}$ \textup{(}all of it\textup{)},
\begin{equation}\label{eq:transfer}
   \dstar(E_v)\ \le\ \frac{\pi^2}{\mu(v)^2}\,\dstar(\Om).
\end{equation}
\end{lemma}

\begin{proof}
On $E_v$ the function $w:=u-v$ solves $-\Delta w=1$, $w=0$ on
$\partial E_v$, so $w$ is the torsion function of $E_v$ and
$T(E_v)=\int_{E_v}(u-v)$. Its super-level sets are $\{w>s\}=\{u>v+s\}$,
hence the level defects of $E_v$ at height $s$ equal $D_H(v+s),D_I(v+s)$.
The scale-invariant form of Fact~\ref{fact:identity} gives
\[
   4\pi\bigl(T(B_{E_v})-T(E_v)\bigr)
   =\int_0^{m-v}\!\bigl(D_H(v+s)+D_I(v+s)\bigr)ds
   =\int_v^m\!\bigl(D_H+D_I\bigr)dv'
   \le\int_0^m\!\bigl(D_H+D_I\bigr)=4\pi\dT,
\]
the last step by positivity of the integrand. Thus
$T(B_{E_v})-T(E_v)\le\dT$. Dividing by $T(B_{E_v})=\mu(v)^2/8\pi$,
\[
   \dstar(E_v)=\frac{T(B_{E_v})-T(E_v)}{T(B_{E_v})}
   \le\frac{\dT}{\mu(v)^2/8\pi}=\frac{8\pi\dT}{\mu(v)^2}
   =\frac{\pi^2}{\mu(v)^2}\cdot\frac{8\dT}{\pi}
   =\frac{\pi^2}{\mu(v)^2}\,\dstar(\Om).\qedhere
\]
\end{proof}

The point is that $\dstar(E_v)\lesssim\dstar(\Om)$ at \emph{every} level (a
tail of a positive-integrand integral), even though the single-level defects
$D_H(v),D_I(v)$ are inflated near the boundary. Removing the small spillover
components (\S\ref{sec:spill}) only decreases the deficit, by
Lemma~\ref{lem:cutdeficit}.

\subsection{Sharp level selection}\label{sec:select}

\begin{lemma}[Good level]\label{lem:level}
There is an absolute $A_0$ \textup{(}one may take $A_0=16\pi$\textup{)} such
that for every $\tau\in(0,\tfrac1{16})$ there is a regular value
$\hat v\in(\tau,2\tau)$ with
\begin{equation}\label{eq:goodlevel}
   D_H(\hat v)+D_I(\hat v)\ \le\ A_0\,\frac{\dT}{\tau},
   \qquad
   \mu_B(\hat v)-\mu(\hat v)\ \le\ A_0\,\frac{\dT}{\tau}.
\end{equation}
\end{lemma}

\begin{proof}
On $I=(\tau,2\tau)$, $\abs I=\tau$. By \eqref{eq:identity},
$\int_I(D_H+D_I)\le4\pi\dT$, so by Chebyshev
$\bigl|\{v\in I:D_H+D_I>A_0\dT/\tau\}\bigr|\le4\pi\dT/(A_0\dT/\tau)
=4\pi\tau/A_0$. By \eqref{eq:profile}, $\int_I(\mu_B-\mu)\le\dT$, so
$\bigl|\{v\in I:\mu_B-\mu>A_0\dT/\tau\}\bigr|\le\tau/A_0$. The non-regular
values have measure $0$ (Fact~\ref{fact:reg}). With $A_0=16\pi$ the union of
the bad sets has measure $\le4\pi\tau/A_0+\tau/A_0=\tfrac14\tau+\tfrac
\tau{16\pi}<\tau=\abs I$, so a good regular $\hat v$ exists.
\end{proof}

\subsection{Closeness}\label{sec:close}

\begin{lemma}[Closeness of the super-level set]\label{lem:close}
At the level $\hat v$ of Lemma~\ref{lem:level},
\begin{equation}\label{eq:close}
   \abs{\Om\triangle\{u>\hat v\}}=\pi-\mu(\hat v)\ \le\ 8\pi\tau+A_0\frac{\dT}{\tau}.
\end{equation}
\end{lemma}

\begin{proof}
$\{u>\hat v\}\subset\Om$, so $\abs{\Om\triangle\{u>\hat v\}}=\abs{\{u\le\hat
v\}}=\pi-\mu(\hat v)$. Split
\[
   \pi-\mu(\hat v)=\bigl(\pi-\mu_B(\hat v)\bigr)+\bigl(\mu_B(\hat v)-\mu(\hat v)\bigr)
   =4\pi\hat v+\bigl(\mu_B(\hat v)-\mu(\hat v)\bigr)
   \le4\pi(2\tau)+A_0\frac{\dT}{\tau},
\]
using $\mu_B(\hat v)=\pi(1-4\hat v)$ (as $\hat v\le2\tau<\tfrac14$) and
\eqref{eq:goodlevel}.
\end{proof}

\subsection{One connected inset: the spillover is higher order}\label{sec:spill}

\begin{lemma}[Spillover]\label{lem:spill}
Let $\Etil$ be the connected component of $\{u>\hat v\}$ of largest measure.
Then
\begin{equation}\label{eq:spill}
   \bigl|\{u>\hat v\}\setminus\Etil\bigr|\ \le\ \frac{D_I(\hat v)^2}{16\pi^3}
   \ \le\ \frac{A_0^2}{16\pi^3}\,\frac{\dT^2}{\tau^2}.
\end{equation}
Consequently $\Etil$ is connected, has no interior holes
\textup{(}Fact~\ref{fact:reg}\textup{)}, and
$\abs{\Om\triangle\Etil}\le 8\pi\tau+A_0\dT/\tau+A_0^2\dT^2/(16\pi^3\tau^2)$.
\end{lemma}

\begin{proof}
Let the components of $\{u>\hat v\}$ have areas $A_1\ge A_2\ge\cdots$, so
$\Etil$ is the $A_1$-component and $\sum A_i=\mu(\hat v)$,
$P(\hat v)=\sum P(\{A_i\})\ge\sum 2\sqrt{\pi A_i}$. Then
\[
   D_I(\hat v)=P^2-4\pi\textstyle\sum A_i
   \ge4\pi\Bigl[\bigl(\sum\sqrt{A_i}\bigr)^2-\sum A_i\Bigr]
   =8\pi\!\!\sum_{i<j}\!\sqrt{A_iA_j}
   \ge8\pi\sqrt{A_1}\!\sum_{i\ge2}\!\sqrt{A_i}
   \ge8\pi\sqrt{A_1}\,\sqrt{\textstyle\sum_{i\ge2}A_i}.
\]
For $\dT\le\delta_\star$ we have $\mu(\hat v)\ge\pi-1\ge\pi/2$ and (since the
spillover is small) $A_1\ge\pi/4$, whence
$\sum_{i\ge2}A_i\le D_I(\hat v)^2/(64\pi^2A_1)\le D_I(\hat v)^2/(16\pi^3)$.
The last estimate uses \eqref{eq:goodlevel}. The remaining claims follow,
adding \eqref{eq:spill} to \eqref{eq:close} and noting
$\abs{\Om\triangle\Etil}=(\pi-\mu(\hat v))+\abs{\{u>\hat v\}\setminus\Etil}$.
\end{proof}

From now on $\Etil$ denotes this connected inset, with
\[
   \mu_E:=\abs\Etil=\mu(\hat v)-\bigl|\{u>\hat v\}\setminus\Etil\bigr|,
   \quad
   \hat r:=\sqrt{\mu_E/\pi},
   \quad
   \diso:=\diso(\Etil),
   \quad
   \Exc:=P(\Etil)-2\pi\hat r .
\]
Since $\Etil$ is a component of $\{u>\hat v\}$,
$P(\Etil)\le P(\hat v)$ and hence
\begin{equation}\label{eq:disoE}
   \diso\ \le\ \frac{D_I(\hat v)}{8\pi\,\mu_E}\ \le\ C\,\frac{\dT}{\tau},
   \qquad
   \Exc=2\pi\hat r\,\diso\le C\,\frac{\dT}{\tau},
\end{equation}
using $P(\hat v)^2=4\pi\mu(\hat v)+D_I(\hat v)$ and
$\sqrt{1+t}-1\le t/2$. Also $\dstar(\Etil)\le\dstar(E_{\hat v})\le
(\pi^2/\mu_E^2)\dstar(\Om)$ by Lemmas~\ref{lem:transfer} and
\ref{lem:cutdeficit} below.

\section{The fold content has no \texorpdfstring{$\sqrt{\,\cdot\,}$}{sqrt}-loss}
\label{sec:fold}

This is the heart of the sharp reduction. We bound the total radial
\emph{fold content} of $\Etil$ — the obstruction to being a graph about $z$ —
by $\Exc+a\eta$, where $a=\abs{\Etil\triangle B_{\hat r}(z)}$ and $\eta$ is
the trapping width. The improvement over the naive $\Exc+Ca$ (which would
carry the quantitative-isoperimetric loss $a\sim\sqrt\diso$) is the factor
$\eta$ multiplying $a$: it comes from the cancellation of
$\int_{\Etil}\abs{x-z}^{-1}$ against $2\pi\hat r$ across the thin annulus.

For a.e.\ ray direction $\theta$ the slice is
$\Etil_\theta=(0,r_1)\cup(r_2,r_3)\cup\cdots\cup(r_{2m},r_{2m+1})$ with
$r_k\in[\rho_i,\rho_e]$ and $M(\theta)=2m(\theta)+1$; the intervals
$(r_{2i},r_{2i+1})$, $i\ge1$, are the \emph{caps}, and $\Etil$ is a radial
graph about $z$ iff $M\equiv1$.

\begin{proposition}[Sharp fold content]\label{prop:fold}
With $z,\rho_i,\rho_e,\eta,\hat r$ as in Fact~\ref{fact:annulus} and
$a=\abs{\Etil\triangle B_{\hat r}(z)}$,
\begin{equation}\label{eq:foldsharp}
   \int_{S^1}\bigl(M(\theta)-1\bigr)\,d\theta
   \ \le\ \frac{1}{\rho_i}\Bigl(\Exc+\frac{a\,\eta}{\rho_i\hat r}\Bigr)
   \ \le\ C\bigl(\Exc+a\eta\bigr).
\end{equation}
\end{proposition}

\begin{proof}
By Fact~\ref{fact:flux}, subtracting the two identities in
\eqref{eq:two-flux},
\[
   \int_{S^1}\!(M-1)\,d\theta
   =\int_{\partial\Etil}\frac{\abs{\nu_r}-\nu_r}{\abs{x-z}}\,d\HH^1
   =2\!\int_{\partial\Etil\cap\{\nu_r<0\}}\!\frac{\abs{\nu_r}}{\abs{x-z}}\,d\HH^1
   \ \le\ \frac{2}{\rho_i}\,F_r,\quad
   F_r:=\!\int_{\{\nu_r<0\}}\!\abs{\nu_r}\,d\HH^1,
\]
since $\abs{x-z}\ge\rho_i$ on $\partial\Etil$. By the third identity in
\eqref{eq:two-flux}, $\int_{\partial\Etil}\nu_r=I_{\Etil}$, while
$\int\nu_r=\int\abs{\nu_r}-2F_r=P_r-2F_r$ with $P_r=\int\abs{\nu_r}$; hence
$2F_r=P_r-I_{\Etil}$.

Now $P_r\le P(\Etil)=2\pi\hat r+\Exc$. For $I_{\Etil}$ we use that
$\abs\Etil=\abs{B_{\hat r}(z)}=\mu_E$, so
$\int(\mathbf1_{\Etil}-\mathbf1_{B_{\hat r}(z)})=0$ and we may subtract the
constant $1/\hat r$:
\[
   I_{\Etil}-2\pi\hat r
   =\int\bigl(\mathbf1_{\Etil}-\mathbf1_{B_{\hat r}(z)}\bigr)
       \frac{1}{\abs{x-z}}
   =\int\bigl(\mathbf1_{\Etil}-\mathbf1_{B_{\hat r}(z)}\bigr)
       \Bigl(\frac{1}{\abs{x-z}}-\frac1{\hat r}\Bigr),
\]
using $\int_{B_{\hat r}(z)}\abs{x-z}^{-1}=2\pi\hat r$. The integrand is
supported on $\Etil\triangle B_{\hat r}(z)$, which lies in the annulus
$\{\rho_i\le\abs{x-z}\le\rho_e\}$ (because $B_{\rho_i}(z)\subset\Etil\cap
B_{\hat r}(z)$ and $\Etil\cup B_{\hat r}(z)\subset B_{\rho_e}(z)$). There
$\bigl|\abs{x-z}^{-1}-\hat r^{-1}\bigr|=\bigl|\hat r-\abs{x-z}\bigr|/
(\hat r\abs{x-z})\le\eta/(\rho_i\hat r)$, since $\hat r,\abs{x-z}\in
[\rho_i,\rho_e]$ and $\rho_e-\rho_i=\eta$. Therefore
\begin{equation}\label{eq:Icancel}
   \bigl|I_{\Etil}-2\pi\hat r\bigr|\ \le\ \frac{\eta}{\rho_i\hat r}\,
   \abs{\Etil\triangle B_{\hat r}(z)}=\frac{a\,\eta}{\rho_i\hat r}.
\end{equation}
Combining, $2F_r=P_r-I_{\Etil}\le(2\pi\hat r+\Exc)-(2\pi\hat r-a\eta/(\rho_i
\hat r))=\Exc+a\eta/(\rho_i\hat r)$, and
$\int(M-1)\le(2/\rho_i)F_r\le(\Exc+a\eta/(\rho_i\hat r))/\rho_i$. As
$\rho_i,\hat r\ge\tfrac12$ for $\dT\le\delta_\star$, this is
$\le C(\Exc+a\eta)$.
\end{proof}

\begin{remark}[Why \eqref{eq:Icancel} is the crux]
The earlier draft bounded $\abs{I_{\Etil}-2\pi\hat r}\le a/\rho_i\le Ca$ by
pulling $\abs{x-z}^{-1}\le\rho_i^{-1}$ out of the integral. That carries the
full $a\sim\sqrt\diso$, i.e.\ the FMP square-root, into the fold content and
ultimately caps the result at $\dT^{1/4}$ closeness. Subtracting the
constant $1/\hat r$ (legitimate because the two indicators have equal mass)
and using that $\Etil\triangle B_{\hat r}$ is confined to the
$\eta$-annulus replaces $a$ by $a\eta$. This single factor $\eta$ is what
makes the reduction sharp.
\end{remark}

\begin{corollary}[Graph defect]\label{cor:gd}
The optimal radial-graph defect of $\Etil$ about $z$ satisfies
\begin{equation}\label{eq:gd}
   \gd(\Etil)\ \le\ \rho_e\,\eta\!\int_{S^1}\!m\,d\theta
   \ =\ \tfrac12\rho_e\,\eta\!\int_{S^1}\!(M-1)\,d\theta
   \ \le\ C\,\eta\,(\Exc+a\eta).
\end{equation}
\end{corollary}

\begin{proof}
Take the core profile $\rho_0(\theta):=r_1(\theta)\in[\rho_i,\rho_e]$ (cut
every cap); then $\{r<\rho_0\}\subset\Etil$ and $\Etil\setminus\{r<\rho_0\}$
is the union of caps. On each ray the caps lie in $[\rho_i,\rho_e]$ and have
total radial length $\le r_{2m+1}-r_1\le\eta$, so the per-ray cap mass is
$\int_{\text{caps}}r\,dr\le\rho_e\,\eta\,\mathbf 1_{\{m\ge1\}}$. Integrating
in $\theta$ and using $\abs{\{m\ge1\}}\le\int m=\tfrac12\int(M-1)$ gives the
first two relations; Proposition~\ref{prop:fold} gives the last.
\end{proof}

\section{Cutting the caps lowers the deficit}
\label{sec:cut}

Manufacturing an honest graph from $\Etil$ means cutting the caps. The
naive fear is that this inflates the torsion deficit by
$\sim\tilde\eta\,\gd$. It does not: removing low-torsion boundary material
shrinks the equal-area comparison ball faster than it loses torsion, so the
\emph{scale-invariant} deficit goes down.

\begin{lemma}[Cutting boundary material decreases $\dstar$]\label{lem:cutdeficit}
Let $B_{\rho_i}(z)\subset G\subset E\subset B_{\rho_e}(z)$ with
$E\setminus G$ contained in the collar $\{\rho_i<\abs{x-z}<\rho_e\}$, and
$\tilde\eta:=\tfrac14(\rho_e^2-\rho_i^2)$. If
\begin{equation}\label{eq:cutcond}
   3\,\tilde\eta\ \le\ \frac{T(E)}{\abs E},
\end{equation}
then $\dstar(G)\le\dstar(E)$.
\end{lemma}

\begin{proof}
Write $V=\abs E$, $g=\abs{E\setminus G}=V-\abs G\in[0,V]$, $T=T(E)$. By
Fact~\ref{fact:torsion-stab}, $T(E)-T(G)\le3\tilde\eta\,g$, i.e.\
$T(G)\ge T-3\tilde\eta\,g$. With $T(B_E)=V^2/8\pi$ and
$T(B_G)=(V-g)^2/8\pi$,
\[
   \dstar(G)-\dstar(E)
   =\frac{T}{T(B_E)}-\frac{T(G)}{T(B_G)}
   =8\pi\Bigl[\frac{T}{V^2}-\frac{T(G)}{(V-g)^2}\Bigr],
\]
so $\dstar(G)\le\dstar(E)$ is equivalent to $T(G)\ge T\,(V-g)^2/V^2=
T(1-g/V)^2$. Expanding $T(1-g/V)^2=T-2Tg/V+Tg^2/V^2$,
\[
   T(G)-T\Bigl(1-\frac gV\Bigr)^2
   \ \ge\ \bigl(T-3\tilde\eta\,g\bigr)-\Bigl(T-\frac{2Tg}{V}+\frac{Tg^2}{V^2}\Bigr)
   = g\Bigl(\frac{2T}{V}-\frac{Tg}{V^2}-3\tilde\eta\Bigr).
\]
Since $g\le V$ we have $\dfrac{2T}{V}-\dfrac{Tg}{V^2}=\dfrac TV\Bigl(2-\dfrac
gV\Bigr)\ge\dfrac TV$, so the bracket is $\ge\dfrac TV-3\tilde\eta\ge0$ by
\eqref{eq:cutcond}. Hence $T(G)\ge T(1-g/V)^2$ and $\dstar(G)\le\dstar(E)$.
\end{proof}

\begin{remark}
Condition \eqref{eq:cutcond} is met with room to spare in our application:
$T(\Etil)/\abs\Etil=(1-\dstar(\Etil))\abs\Etil/(8\pi)\ge(1-\tfrac12)
(\pi/2)/(8\pi)=\tfrac1{32}$, while $\tilde\eta\le\eta\le C(\dT/\tau)^{1/4}
\to0$. So $3\tilde\eta\le\tfrac1{32}\le T/\abs E$ for $\dT\le\delta_\star$.
Lemma~\ref{lem:cutdeficit} also justifies, a posteriori, that passing from
$\{u>\hat v\}$ to its main component $\Etil$ (cutting the spillover, which
sits in the same annulus) does not increase $\dstar$, completing the chain
$\dstar(\Etil)\le\dstar(E_{\hat v})\le(\pi^2/\mu_E^2)\dstar(\Om)$.
\end{remark}

\section{The graph inset and the conclusion}
\label{sec:concl}

\begin{theorem}[Sharp reduction to a graph inset]\label{thm:reduction}
There is $\delta_\star>0$ absolute so that for $\Om\subset\R^2$,
$\abs\Om=\pi$, $\dT\le\delta_\star$, with $\tau:=\sqrt\dT$ there exist a
point $z$ and a continuous radial graph
$G=\{z+r\omega:r<\rho(\omega)\}$, $\rho\in C^0(S^1)$, with
\begin{enumerate}[label=\textup{(\roman*)},leftmargin=2.4em]
\item \textup{(small amplitude)}
   $\norm{\rho/\hat r_G-1}_{L^\infty}\le C\,\dT^{1/8}\le\eta_G$,
   $\hat r_G=\sqrt{\abs G/\pi}$;
\item \textup{(controlled deficit)} $\dstar(G)\ \le\ C\,\dstar(\Om)$;
\item \textup{(closeness)} $\abs{\Om\triangle G}\ \le\ C\,\sqrt\dT$.
\end{enumerate}
Every constant is absolute, and the construction uses only $D_H$ and $D_I$
(via Facts~\ref{fact:identity}--\ref{fact:annulus} and \ref{fact:flux}); it
uses no selection principle, no transport, and no Magnanini--Poggesi
estimate.
\end{theorem}

\begin{proof}
Fix $\tau=\sqrt\dT$ and let $\hat v\in(\tau,2\tau)$, $\Etil$ be the level
and connected inset of \S\ref{sec:inset}. We collect the scales. By
\eqref{eq:disoE}, $\diso=\diso(\Etil)\le C\dT/\tau=C\sqrt\dT$ and
$\Exc\le C\sqrt\dT$. By Fact~\ref{fact:fmp},
$a=\abs{\Etil\triangle B_{\hat r}(z)}\le C\sqrt{\diso}\le C\dT^{1/4}$. By
Fact~\ref{fact:annulus},
\begin{equation}\label{eq:etascale}
   \eta=\rho_e-\rho_i\le C(\Exc+a)^{1/2}\le C\,a^{1/2}\le C\,\dT^{1/8},
\end{equation}
since $a\sim\dT^{1/4}$ dominates $\Exc\sim\dT^{1/2}$. By
Corollary~\ref{cor:gd},
\begin{equation}\label{eq:gdscale}
   \gd(\Etil)\le C\,\eta(\Exc+a\eta)
   \le C\bigl(\dT^{1/8}\!\cdot\dT^{1/2}+\dT^{1/4}\!\cdot\dT^{1/4}\bigr)
   =C\bigl(\dT^{5/8}+\dT^{1/2}\bigr)\le C\sqrt\dT .
\end{equation}

\emph{The graph.} Let $\rho_0=r_1\in BV(S^1;[\rho_i,\rho_e])$ be the core
profile (Corollary~\ref{cor:gd}), $G_0:=\{r<\rho_0\}\subset\Etil$, so
$\abs{\Etil\triangle G_0}=\gd(\Etil)\le C\sqrt\dT$ and
$B_{\rho_i}(z)\subset G_0$. Mollify $\rho_0$ at angular scale
$\sigma=\dT$ to a continuous $\rho\in C^0(S^1;[\rho_i,\rho_e])$; since
$\mathrm{TV}(\rho_0)\le P(\Etil)/\rho_i\le C$, this costs
$\abs{G_0\triangle G}\le\sigma\,\mathrm{TV}(\rho_0)\,\rho_e\le C\dT$, where
$G=\{r<\rho\}$. Thus $\abs{\Etil\triangle G}\le\gd(\Etil)+C\dT\le C\sqrt\dT$
and, with Lemma~\ref{lem:spill} and \eqref{eq:close},
\[
   \abs{\Om\triangle G}\le\abs{\Om\triangle\Etil}+\abs{\Etil\triangle G}
   \le\Bigl(8\pi\tau+A_0\tfrac\dT\tau+C\tfrac{\dT^2}{\tau^2}\Bigr)+C\sqrt\dT
   \le C\sqrt\dT,
\]
which is (iii). The amplitude obeys $\norm{\rho-\hat r}_\infty\le\rho_e-
\rho_i=\eta\le C\dT^{1/8}$, and $\hat r_G=\sqrt{\abs G/\pi}=\hat r+O(\sqrt
\dT)$, so $\norm{\rho/\hat r_G-1}_\infty\le C\dT^{1/8}\le\eta_G$ for
$\dT\le\delta_\star$: this is (i).

\emph{Deficit.} By Lemma~\ref{lem:cutdeficit} applied to $G_0\subset\Etil$
(its symmetric difference is the caps, in the collar; condition
\eqref{eq:cutcond} holds by the Remark), $\dstar(G_0)\le\dstar(\Etil)$. The
mollification changes area and torsion by $O(\dT)$ in the collar, so by
Fact~\ref{fact:torsion-stab} and \eqref{eq:dstar},
\[
   \dstar(G)\le\dstar(G_0)+C\bigl(\tilde\eta\abs{G\triangle G_0}
   +\abs{\,\abs G-\abs{G_0}\,}\bigr)\le\dstar(\Etil)+C\dT.
\]
Finally $\dstar(\Etil)\le(\pi^2/\mu_E^2)\dstar(\Om)\le(1+C\sqrt\dT)\dstar(\Om)$
by Lemma~\ref{lem:transfer} (with $\mu_E=\pi-O(\sqrt\dT)$) and
Lemma~\ref{lem:cutdeficit}. Hence $\dstar(G)\le C\dstar(\Om)$, which is
(ii), since $\dT\le C\dstar(\Om)$.
\end{proof}

\begin{theorem}[Sharp quantitative Saint--Venant, selection-free]
\label{thm:main}
There is an absolute constant $C$ such that every open $\Om\subset\R^2$ with
$\abs\Om=\pi$ satisfies
\begin{equation}\label{eq:main}
   \AF(\Om)^2\ \le\ C\,\dT(\Om).
\end{equation}
\end{theorem}

\begin{proof}
For $\dT\ge\delta_\star$ the bound is trivial
($\AF\le2$, so $\AF^2\le4\le(4/\delta_\star)\dT$). For $\dT\le\delta_\star$,
let $G$ be the graph of Theorem~\ref{thm:reduction}, centred at $z$ with
barycentre within $O(\eta)\le O(\dT^{1/8})$ of $z$; recentre by a
translation of size $O(\dT^{1/8})\le\eta_G$ so that $G$ has barycentre $z$
(this changes $\abs{\Om\triangle G}$ by $O(\dT^{1/8})\cdot P(G)=O(\dT^{1/8})$
— absorbed into a slightly larger constant in (iii) after re-optimising
$\tau$; equivalently apply (G) in the barycentre-absorbed form, which costs
only $\dstar(G)$). By the radial-graph theorem (Fact~\ref{fact:graph}) and
(ii),
\[
   \AF(G)^2\ \le\ C_G\,\dstar(G)\ \le\ C\,\dstar(\Om)\ =\ C\,\tfrac8\pi\dT .
\]
Let $B$ be the disc of area $\pi$ centred at the barycentre of $G$. Then,
using $\AF(G)=\abs{G\triangle B_G}/\abs G$ for the equal-area disc $B_G$
and $\abs{B\triangle B_G}=\bigl|\,\pi-\abs G\,\bigr|\le\abs{\Om\triangle
G}\le C\sqrt\dT$,
\[
   \AF(\Om)\le\frac{\abs{\Om\triangle B}}{\pi}
   \le\frac{\abs{\Om\triangle G}+\abs{G\triangle B_G}+\abs{B_G\triangle B}}{\pi}
   \le\frac1\pi\bigl(C\sqrt\dT+\AF(G)\abs G+C\sqrt\dT\bigr)
   \le C\sqrt\dT .
\]
Squaring gives \eqref{eq:main}.
\end{proof}

\begin{remark}[What did the work, and what was \emph{not} used]
The reduction Theorem~\ref{thm:reduction} draws only on the two level
defects:
\begin{itemize}[leftmargin=1.4em]
\item $D_H+D_I$ through the identity \eqref{eq:identity} (level selection,
deficit transfer) and through $-\mu'\ge4\pi$;
\item $D_I$ through the isoperimetric/annulus inputs
(Facts~\ref{fact:annulus},~\ref{fact:fmp}) and the multi-component spillover
bound.
\end{itemize}
The two sharp moves are \eqref{eq:Icancel} (fold content without the FMP
square-root) and Lemma~\ref{lem:cutdeficit} (cutting lowers $\dstar$). With
$\tau=\sqrt\dT$, closeness $8\pi\tau$ and graph defect $C\dT/\tau$ are both
$O(\sqrt\dT)$; this balance is the source of the exponent. No moving-ball or
kinetic transport, no Magnanini--Poggesi inradius/outradius estimate, and no
Brasco--De~Philippis--Velichkov selection principle enter at any point. The
perturbative finish is exactly the radial-graph theorem (G).
\end{remark}

\begin{remark}[Where the previous obstructions dissolved]
Two walls had blocked the sharp bound. (1) At the boundary scale the
single-level defect inflates, $D_I(\hat v)\sim\dT/\tau\sim\sqrt\dT$, so the
inset is only $\diso$-round to $O(\dT^{1/4})$ in $L^1$ (FMP); but the
graph theorem needs only the \emph{torsion} deficit $\dstar(\Etil)\lesssim
\dT$ (a tail integral), and only an \emph{absolute} amplitude
$\eta\sim\dT^{1/8}\le\eta_G$ — neither needs the $L^1$-distance. Closeness is
carried by the \emph{level} $\hat v\sim\sqrt\dT$, not by $\eta$. (2) Cutting
the unavoidable boundary-layer folds seemed to cost
$\tilde\eta\,\gd\sim\dT^{5/8}$ in deficit; but the equal-area ball shrinks by
$\gd/4$, dominating the torsion loss, so $\dstar$ \emph{decreases}
(Lemma~\ref{lem:cutdeficit}). Decoupling closeness (level) from amplitude
(annulus), and the correct sign in (2), are what close the gap.
\end{remark}

\end{document}
